\pdfoutput=1
%% Author: PGL  Porta Mana
%% Created: 2015-05-01T20:53:34+0200
%% Last-Updated: 2018-12-31T09:30:13+0100
%%%%%%%%%%%%%%%%%%%%%%%%%%%%%%%%%%%%%%%%%%%%%%%%%%%%%%%%%%%%%%%%%%%%%%
% Report-no: ***
\newif\ifarxiv
\arxivfalse
\ifarxiv\pdfmapfile{+classico.map}\fi
\newif\ifafour
\afourfalse % true = A4, false = A5
\newif\iftypodisclaim % typographical disclaim on the side
\typodisclaimtrue
\newif\ifpublic
\publictrue% true = for publication, false = personal notes
\newcommand*{\memfontfamily}{zplx}
\newcommand*{\memfontpack}{newpxtext}
\documentclass[\ifafour a4paper,12pt,\else a5paper,10pt,\fi%extrafontsizes,%
onecolumn,oneside,article,%french,italian,german,swedish,latin,
british%
]{memoir}
\newcommand*{\updated}{\today}
\newcommand*{\firstdraft}{6 August 2018}
\newcommand*{\firstpublished}{\today}
\newcommand*{\propertitle}{Generalization error and overtraining\\{\Large a
  geometric and Bayesian understanding}}
\newcommand*{\pdftitle}{Generalization error and overtraining: a geometric and Bayesian understanding}
\newcommand*{\headtitle}{Bayesian generalization error and overtraining}
\newcommand*{\pdfauthor}{P.G.L.  Porta Mana, I. A. Davidovich}
\newcommand*{\headauthor}{\ifpublic Porta Mana \amp\ Davidovich%
\else Luca \amp\ Andr\'es\fi}
\newcommand*{\reporthead}{}
%%%%%%%%%%%%%%%%%%%%%%%%%%%%%%%%%%%%%%%%%%%%%%%%%%%%%%%%%%%%%%%%%%%%%%%%%%%%
%%%%%%%%%%%%%%%%%%%%%%%%%%%%%%%%%%%%%%%%%%%%%%%%%%%%%%%%%%%%%%%%%%%%%%%%%%%%
%\usepackage{pifont}
%\usepackage{fontawesome}
\usepackage[T1]{fontenc} 
\input{glyphtounicode} \pdfgentounicode=1
\usepackage[utf8]{inputenx}
%\usepackage{newunicodechar}
% \newunicodechar{Ĕ}{\u{E}}
% \newunicodechar{ĕ}{\u{e}}
% \newunicodechar{Ĭ}{\u{I}}
% \newunicodechar{ĭ}{\u{\i}}
% \newunicodechar{Ŏ}{\u{O}}
% \newunicodechar{ŏ}{\u{o}}
% \newunicodechar{Ŭ}{\u{U}}
% \newunicodechar{ŭ}{\u{u}}
% \newunicodechar{Ā}{\=A}
% \newunicodechar{ā}{\=a}
% \newunicodechar{Ē}{\=E}
% \newunicodechar{ē}{\=e}
% \newunicodechar{Ī}{\=I}
% \newunicodechar{ī}{\={\i}}
% \newunicodechar{Ō}{\=O}
% \newunicodechar{ō}{\=o}
% \newunicodechar{Ū}{\=U}
% \newunicodechar{ū}{\=u}
% \newunicodechar{Ȳ}{\=Y}
% \newunicodechar{ȳ}{\=y}

\newcommand*{\bmmax}{0} % reduce number of bold fonts, before bm
\newcommand*{\hmmax}{0} % reduce number of heavy fonts, before bm
\usepackage{textcomp}
\usepackage[normalem]{ulem}
% \makeatletter
% \def\ssout{\bgroup \ULdepth=-.35ex%\UL@setULdepth
%  \markoverwith{\lower\ULdepth\hbox
%    {\kern-.03em\vbox{\hrule width.2em\kern1.2\p@\hrule}\kern-.03em}}%
%  \ULon}
% \makeatother
\usepackage{amsmath}
\usepackage{mathtools}
\addtolength{\jot}{\jot} % increase spacing in multiline formulae
\usepackage{empheq}% automatically calls amsmath and mathtools
\newcommand*{\widefbox}[1]{\fbox{\hspace{1em}#1\hspace{1em}}}
\setlength{\multlinegap}{0pt}
%\usepackage{fancybox}
%\usepackage{framed}
% \usepackage[misc]{ifsym} % for dice
% \newcommand*{\diceone}{{\scriptsize\Cube{1}}}
\usepackage{amssymb}
\usepackage{amsxtra}

\usepackage[main=british,french,italian,german,swedish,latin,esperanto]{babel}\selectlanguage{british}
\newcommand*{\langfrench}{\foreignlanguage{french}}
\newcommand*{\langgerman}{\foreignlanguage{german}}
\newcommand*{\langitalian}{\foreignlanguage{italian}}
\newcommand*{\langswedish}{\foreignlanguage{swedish}}
\newcommand*{\langlatin}{\foreignlanguage{latin}}
\newcommand*{\langnohyph}{\foreignlanguage{nohyphenation}}

\usepackage[autostyle=false,autopunct=false,english=british]{csquotes}
\setquotestyle{british}

\usepackage{amsthm}
\newcommand*{\QED}{\textsc{q.e.d.}}
\renewcommand*{\qedsymbol}{\QED}
\theoremstyle{remark}
\newtheorem{note}{Note}
\newtheorem*{remark}{Note}
\newtheoremstyle{innote}{\parsep}{\parsep}{\footnotesize}{}{}{}{0pt}{}
\theoremstyle{innote}
\newtheorem*{innote}{}


\usepackage[shortlabels,inline]{enumitem}
\SetEnumitemKey{para}{itemindent=\parindent,leftmargin=0pt,listparindent=\parindent,parsep=0pt,itemsep=\topsep}
% \begin{asparaenum} = \begin{enumerate}[para]
% \begin{inparaenum} = \begin{enumerate*}
\setlist[enumerate,2]{label=\alph*.}
\setlist[enumerate]{leftmargin=\parindent}
\setlist[itemize]{leftmargin=\parindent}
\setlist[description]{leftmargin=\parindent}

\usepackage[babel,theoremfont]{newpxtext}
\usepackage[bigdelims,nosymbolsc%,smallerops % probably arXiv doesn't have it
]{newpxmath}
\useosf\linespread{1.083}
%% smaller operators for old version of newpxmath
\makeatletter
\def\re@DeclareMathSymbol#1#2#3#4{%
    \let#1=\undefined
    \DeclareMathSymbol{#1}{#2}{#3}{#4}}
%\re@DeclareMathSymbol{\bigsqcupop}{\mathop}{largesymbols}{"46}
%\re@DeclareMathSymbol{\bigodotop}{\mathop}{largesymbols}{"4A}
\re@DeclareMathSymbol{\bigoplusop}{\mathop}{largesymbols}{"4C}
\re@DeclareMathSymbol{\bigotimesop}{\mathop}{largesymbols}{"4E}
\re@DeclareMathSymbol{\sumop}{\mathop}{largesymbols}{"50}
\re@DeclareMathSymbol{\prodop}{\mathop}{largesymbols}{"51}
\re@DeclareMathSymbol{\bigcupop}{\mathop}{largesymbols}{"53}
\re@DeclareMathSymbol{\bigcapop}{\mathop}{largesymbols}{"54}
%\re@DeclareMathSymbol{\biguplusop}{\mathop}{largesymbols}{"55}
\re@DeclareMathSymbol{\bigwedgeop}{\mathop}{largesymbols}{"56}
\re@DeclareMathSymbol{\bigveeop}{\mathop}{largesymbols}{"57}
%\re@DeclareMathSymbol{\bigcupdotop}{\mathop}{largesymbols}{"DF}
%\re@DeclareMathSymbol{\bigcapplusop}{\mathop}{largesymbolsPXA}{"00}
%\re@DeclareMathSymbol{\bigsqcupplusop}{\mathop}{largesymbolsPXA}{"02}
%\re@DeclareMathSymbol{\bigsqcapplusop}{\mathop}{largesymbolsPXA}{"04}
%\re@DeclareMathSymbol{\bigsqcapop}{\mathop}{largesymbolsPXA}{"06}
\re@DeclareMathSymbol{\bigtimesop}{\mathop}{largesymbolsPXA}{"10}
%\re@DeclareMathSymbol{\coprodop}{\mathop}{largesymbols}{"60}
%\re@DeclareMathSymbol{\varprod}{\mathop}{largesymbolsPXA}{16}
\makeatother


%% With euler font cursive for Greek letters - the [1] means 100% scaling
\DeclareFontFamily{U}{egreek}{\skewchar\font'177}%
\DeclareFontShape{U}{egreek}{m}{n}{<-6>s*[1]eurm5 <6-8>s*[1]eurm7 <8->s*[1]eurm10}{}%
\DeclareFontShape{U}{egreek}{m}{it}{<->s*[1]eurmo10}{}%
\DeclareFontShape{U}{egreek}{b}{n}{<-6>s*[1]eurb5 <6-8>s*[1]eurb7 <8->s*[1]eurb10}{}%
\DeclareFontShape{U}{egreek}{b}{it}{<->s*[1]eurbo10}{}%
\DeclareSymbolFont{egreeki}{U}{egreek}{m}{it}%
\SetSymbolFont{egreeki}{bold}{U}{egreek}{b}{it}% from the amsfonts package
\DeclareSymbolFont{egreekr}{U}{egreek}{m}{n}%
\SetSymbolFont{egreekr}{bold}{U}{egreek}{b}{n}% from the amsfonts package
% Take also \sum, \prod, \coprod symbols from Euler fonts
\DeclareFontFamily{U}{egreekx}{\skewchar\font'177}
\DeclareFontShape{U}{egreekx}{m}{n}{%
       <-7.5>s*[0.9]euex7%
    <7.5-8.5>s*[0.9]euex8%
    <8.5-9.5>s*[0.9]euex9%
    <9.5->s*[0.9]euex10%
}{}
\DeclareSymbolFont{egreekx}{U}{egreekx}{m}{n}
\DeclareMathSymbol{\sumop}{\mathop}{egreekx}{"50}
\DeclareMathSymbol{\prodop}{\mathop}{egreekx}{"51}
\DeclareMathSymbol{\coprodop}{\mathop}{egreekx}{"60}
\makeatletter
\def\sum{\DOTSI\sumop\slimits@}
\def\prod{\DOTSI\prodop\slimits@}
\def\coprod{\DOTSI\coprodop\slimits@}
\makeatother
\ifarxiv\else\let\varpartial\undefined
\let\partialup\undefined
\let\alpha\undefined
\let\beta\undefined
\let\gamma\undefined
\let\delta\undefined
\let\epsilon\undefined
\let\zeta\undefined
\let\eta\undefined
\let\theta\undefined
\let\iota\undefined
\let\kappa\undefined
\let\lambda\undefined
\let\mu\undefined
\let\nu\undefined
\let\xi\undefined
\let\omicron\undefined
\let\pi\undefined
\let\rho\undefined
\let\sigma\undefined
\let\tau\undefined
\let\upsilon\undefined
\let\phi\undefined
\let\chi\undefined
\let\psi\undefined
\let\omega\undefined
\let\varepsilon\undefined
\let\vartheta\undefined
\let\varpi\undefined
\let\varrho\undefined 
\let\varsigma\undefined
\let\varkappa\undefined
\let\varphi\undefined
%
\let\varAlpha\undefined
\let\varBeta\undefined
\let\varGamma\undefined
\let\varDelta\undefined
\let\varEpsilon\undefined
\let\varZeta\undefined
\let\varEta\undefined
\let\varTheta\undefined
\let\varIota\undefined
\let\varKappa\undefined
\let\varLambda\undefined
\let\varMu\undefined
\let\varNu\undefined
\let\varXi\undefined
\let\varOmicron\undefined
\let\varPi\undefined
\let\varRho\undefined
\let\varSigma\undefined
\let\varTau\undefined
\let\varUpsilon\undefined
\let\varPhi\undefined
\let\varChi\undefined
\let\varPsi\undefined
\let\varOmega\undefined
%
\let\Alpha\undefined
\let\Beta\undefined
\let\Gamma\undefined
\let\Delta\undefined
\let\Epsilon\undefined
\let\Zeta\undefined
\let\Eta\undefined
\let\Theta\undefined
\let\Iota\undefined
\let\Kappa\undefined
\let\Lambda\undefined
\let\Mu\undefined
\let\Nu\undefined
\let\Xi\undefined
\let\Omicron\undefined
\let\Pi\undefined
\let\Rho\undefined
\let\Sigma\undefined
\let\Tau\undefined
\let\Upsilon\undefined
\let\Phi\undefined
\let\Chi\undefined
\let\Psi\undefined
\let\Omega\undefined
%
\let\alphaup\undefined
\let\betaup\undefined
\let\gammaup\undefined
\let\deltaup\undefined
\let\epsilonup\undefined
\let\zetaup\undefined
\let\etaup\undefined
\let\thetaup\undefined
\let\iotaup\undefined
\let\kappaup\undefined
\let\lambdaup\undefined
\let\muup\undefined
\let\nuup\undefined
\let\xiup\undefined
\let\omicronup\undefined
\let\piup\undefined
\let\rhoup\undefined
\let\sigmaup\undefined
\let\tauup\undefined
\let\upsilonup\undefined
\let\phiup\undefined
\let\chiup\undefined
\let\psiup\undefined
\let\omegaup\undefined
\let\varepsilonup\undefined
\let\varthetaup\undefined
\let\varpiup\undefined
\let\varrhoup\undefined
\let\varsigmaup\undefined
\let\varkappaup\undefined
\let\varphiup\undefined
\fi% make sure no CMF greek letters sneak in
% Greek letters not usually given in LaTeX. Comment the unneeded ones
% \DeclareMathSymbol{\varpartial}{\mathalpha}{egreeki}{"40}
 \DeclareMathSymbol{\partialup}{\mathalpha}{egreekr}{"40}
% \DeclareMathSymbol{\alpha}{\mathalpha}{egreeki}{"0B}
% \DeclareMathSymbol{\beta}{\mathalpha}{egreeki}{"0C}
% \DeclareMathSymbol{\gamma}{\mathalpha}{egreeki}{"0D}
% \DeclareMathSymbol{\delta}{\mathalpha}{egreeki}{"0E}
 \DeclareMathSymbol{\epsilon}{\mathalpha}{egreeki}{"0F}
% \DeclareMathSymbol{\zeta}{\mathalpha}{egreeki}{"10}
% \DeclareMathSymbol{\eta}{\mathalpha}{egreeki}{"11}
% \DeclareMathSymbol{\theta}{\mathalpha}{egreeki}{"12}
% \DeclareMathSymbol{\iota}{\mathalpha}{egreeki}{"13}
 \DeclareMathSymbol{\kappa}{\mathalpha}{egreeki}{"14}
 \DeclareMathSymbol{\lambda}{\mathalpha}{egreeki}{"15}
% \DeclareMathSymbol{\mu}{\mathalpha}{egreeki}{"16}
 \DeclareMathSymbol{\nu}{\mathalpha}{egreeki}{"17}
% \DeclareMathSymbol{\xi}{\mathalpha}{egreeki}{"18}
% \DeclareMathSymbol{\omicron}{\mathalpha}{egreeki}{"6F}
% \DeclareMathSymbol{\pi}{\mathalpha}{egreeki}{"19}
% \DeclareMathSymbol{\rho}{\mathalpha}{egreeki}{"1A}
% \DeclareMathSymbol{\sigma}{\mathalpha}{egreeki}{"1B}
% \DeclareMathSymbol{\tau}{\mathalpha}{egreeki}{"1C}
% \DeclareMathSymbol{\upsilon}{\mathalpha}{egreeki}{"1D}
% \DeclareMathSymbol{\phi}{\mathalpha}{egreeki}{"1E}
% \DeclareMathSymbol{\chi}{\mathalpha}{egreeki}{"1F}
% \DeclareMathSymbol{\psi}{\mathalpha}{egreeki}{"20}
% \DeclareMathSymbol{\omega}{\mathalpha}{egreeki}{"21}
% \DeclareMathSymbol{\varepsilon}{\mathalpha}{egreeki}{"22}
% \DeclareMathSymbol{\vartheta}{\mathalpha}{egreeki}{"23}
% \DeclareMathSymbol{\varpi}{\mathalpha}{egreeki}{"24}
% \let\varrho\rho 
% \let\varsigma\sigma
 \let\varkappa\kappa
% \DeclareMathSymbol{\varphi}{\mathalpha}{egreeki}{"27}
% %
% \DeclareMathSymbol{\varAlpha}{\mathalpha}{egreeki}{"41}
% \DeclareMathSymbol{\varBeta}{\mathalpha}{egreeki}{"42}
% \DeclareMathSymbol{\varGamma}{\mathalpha}{egreeki}{"00}
 \DeclareMathSymbol{\varDelta}{\mathalpha}{egreeki}{"01}
 \DeclareMathSymbol{\varEpsilon}{\mathalpha}{egreeki}{"45}
% \DeclareMathSymbol{\varZeta}{\mathalpha}{egreeki}{"5A}
% \DeclareMathSymbol{\varEta}{\mathalpha}{egreeki}{"48}
% \DeclareMathSymbol{\varTheta}{\mathalpha}{egreeki}{"02}
 \DeclareMathSymbol{\varIota}{\mathalpha}{egreeki}{"49}
% \DeclareMathSymbol{\varKappa}{\mathalpha}{egreeki}{"4B}
% \DeclareMathSymbol{\varLambda}{\mathalpha}{egreeki}{"03}
% \DeclareMathSymbol{\varMu}{\mathalpha}{egreeki}{"4D}
% \DeclareMathSymbol{\varNu}{\mathalpha}{egreeki}{"4E}
% \DeclareMathSymbol{\varXi}{\mathalpha}{egreeki}{"04}
% \DeclareMathSymbol{\varOmicron}{\mathalpha}{egreeki}{"4F}
% \DeclareMathSymbol{\varPi}{\mathalpha}{egreeki}{"05}
% \DeclareMathSymbol{\varRho}{\mathalpha}{egreeki}{"50}
% \DeclareMathSymbol{\varSigma}{\mathalpha}{egreeki}{"06}
% \DeclareMathSymbol{\varTau}{\mathalpha}{egreeki}{"54}
% \DeclareMathSymbol{\varUpsilon}{\mathalpha}{egreeki}{"07}
% \DeclareMathSymbol{\varPhi}{\mathalpha}{egreeki}{"08}
% \DeclareMathSymbol{\varChi}{\mathalpha}{egreeki}{"58}
% \DeclareMathSymbol{\varPsi}{\mathalpha}{egreeki}{"09}
% \DeclareMathSymbol{\varOmega}{\mathalpha}{egreeki}{"0A} 
% %
% \DeclareMathSymbol{\Alpha}{\mathalpha}{egreekr}{"41}
% \DeclareMathSymbol{\Beta}{\mathalpha}{egreekr}{"42}
% \DeclareMathSymbol{\Gamma}{\mathalpha}{egreekr}{"00}
% \DeclareMathSymbol{\Delta}{\mathalpha}{egreekr}{"01}
% \DeclareMathSymbol{\Epsilon}{\mathalpha}{egreekr}{"45}
% \DeclareMathSymbol{\Zeta}{\mathalpha}{egreekr}{"5A}
% \DeclareMathSymbol{\Eta}{\mathalpha}{egreekr}{"48}
% \DeclareMathSymbol{\Theta}{\mathalpha}{egreekr}{"02}
% \DeclareMathSymbol{\Iota}{\mathalpha}{egreekr}{"49}
% \DeclareMathSymbol{\Kappa}{\mathalpha}{egreekr}{"4B}
% \DeclareMathSymbol{\Lambda}{\mathalpha}{egreekr}{"03}
% \DeclareMathSymbol{\Mu}{\mathalpha}{egreekr}{"4D}
% \DeclareMathSymbol{\Nu}{\mathalpha}{egreekr}{"4E}
% \DeclareMathSymbol{\Xi}{\mathalpha}{egreekr}{"04}
% \DeclareMathSymbol{\Omicron}{\mathalpha}{egreekr}{"4F}
% \DeclareMathSymbol{\Pi}{\mathalpha}{egreekr}{"05}
% \DeclareMathSymbol{\Rho}{\mathalpha}{egreekr}{"50}
% \DeclareMathSymbol{\Sigma}{\mathalpha}{egreekr}{"06}
% \DeclareMathSymbol{\Tau}{\mathalpha}{egreekr}{"54}
% \DeclareMathSymbol{\Upsilon}{\mathalpha}{egreekr}{"07}
% \DeclareMathSymbol{\Phi}{\mathalpha}{egreekr}{"08}
% \DeclareMathSymbol{\Chi}{\mathalpha}{egreekr}{"58}
% \DeclareMathSymbol{\Psi}{\mathalpha}{egreekr}{"09}
% \DeclareMathSymbol{\Omega}{\mathalpha}{egreekr}{"0A}
% %
% \DeclareMathSymbol{\alphaup}{\mathalpha}{egreekr}{"0B}
% \DeclareMathSymbol{\betaup}{\mathalpha}{egreekr}{"0C}
% \DeclareMathSymbol{\gammaup}{\mathalpha}{egreekr}{"0D}
 \DeclareMathSymbol{\deltaup}{\mathalpha}{egreekr}{"0E}
% \DeclareMathSymbol{\epsilonup}{\mathalpha}{egreekr}{"0F}
% \DeclareMathSymbol{\zetaup}{\mathalpha}{egreekr}{"10}
% \DeclareMathSymbol{\etaup}{\mathalpha}{egreekr}{"11}
% \DeclareMathSymbol{\thetaup}{\mathalpha}{egreekr}{"12}
% \DeclareMathSymbol{\iotaup}{\mathalpha}{egreekr}{"13}
% \DeclareMathSymbol{\kappaup}{\mathalpha}{egreekr}{"14}
% \DeclareMathSymbol{\lambdaup}{\mathalpha}{egreekr}{"15}
% \DeclareMathSymbol{\muup}{\mathalpha}{egreekr}{"16}
% \DeclareMathSymbol{\nuup}{\mathalpha}{egreekr}{"17}
% \DeclareMathSymbol{\xiup}{\mathalpha}{egreekr}{"18}
% \DeclareMathSymbol{\omicronup}{\mathalpha}{egreekr}{"6F}
  \DeclareMathSymbol{\piup}{\mathalpha}{egreekr}{"19}
% \DeclareMathSymbol{\rhoup}{\mathalpha}{egreekr}{"1A}
% \DeclareMathSymbol{\sigmaup}{\mathalpha}{egreekr}{"1B}
% \DeclareMathSymbol{\tauup}{\mathalpha}{egreekr}{"1C}
% \DeclareMathSymbol{\upsilonup}{\mathalpha}{egreekr}{"1D}
% \DeclareMathSymbol{\phiup}{\mathalpha}{egreekr}{"1E}
% \DeclareMathSymbol{\chiup}{\mathalpha}{egreekr}{"1F}
% \DeclareMathSymbol{\psiup}{\mathalpha}{egreekr}{"20}
% \DeclareMathSymbol{\omegaup}{\mathalpha}{egreekr}{"21}
% \DeclareMathSymbol{\varepsilonup}{\mathalpha}{egreekr}{"22}
% \DeclareMathSymbol{\varthetaup}{\mathalpha}{egreekr}{"23}
% \DeclareMathSymbol{\varpiup}{\mathalpha}{egreekr}{"24}
% \let\varrhoup\rhoup 
% \let\varsigmaup\sigmaup
% \let\varkappaup\kappaup
% \DeclareMathSymbol{\varphiup}{\mathalpha}{egreekr}{"27}

% Optima as sans-serif font
%\usepackage%[scaled=0.9]%
%{classico}
\renewcommand\sfdefault{uop}
\DeclareMathAlphabet{\mathsf}  {T1}{\sfdefault}{m}{sl}
\SetMathAlphabet{\mathsf}{bold}{T1}{\sfdefault}{b}{sl}
\newcommand*{\mathte}[1]{\textbf{\textit{\textsf{#1}}}}
% Upright sans-serif math alphabet
% \DeclareMathAlphabet{\mathsu}  {T1}{\sfdefault}{m}{n}
% \SetMathAlphabet{\mathsu}{bold}{T1}{\sfdefault}{b}{n}

% DejaVu Mono as typewriter text
\usepackage[scaled=0.84]{DejaVuSansMono}


\usepackage{mathdots}

\usepackage[usenames]{xcolor}
% Tol (2012) colour-blind-, print-, screen-friendly colours, alternative scheme; Munsell terminology
\definecolor{mypurpleblue}{RGB}{68,119,170}
\definecolor{myblue}{RGB}{102,204,238}
\definecolor{mygreen}{RGB}{34,136,51}
\definecolor{myyellow}{RGB}{204,187,68}
\definecolor{myred}{RGB}{238,102,119}
\definecolor{myredpurple}{RGB}{170,51,119}
\definecolor{mygrey}{RGB}{187,187,187}

% Tol (2012) colour-blind-, print-, screen-friendly colours; Munsell terminology
% \definecolor{lbpurple}{RGB}{51,34,136}
% \definecolor{lblue}{RGB}{136,204,238}
% \definecolor{lbgreen}{RGB}{68,170,153}
% \definecolor{lgreen}{RGB}{17,119,51}
% \definecolor{lgyellow}{RGB}{153,153,51}
% \definecolor{lyellow}{RGB}{221,204,119}
% \definecolor{lred}{RGB}{204,102,119}
% \definecolor{lpred}{RGB}{136,34,85}
% \definecolor{lrpurple}{RGB}{170,68,153}
 \definecolor{lgrey}{RGB}{221,221,221}
%\newcommand*\mycolourbox[1]{%
%\colorbox{mygrey}{\hspace{1em}#1\hspace{1em}}}
\colorlet{shadecolor}{lgrey}

\usepackage{bm}
\usepackage{microtype}

\usepackage[backend=biber,mcite,%subentry,
citestyle=authoryear-comp,bibstyle=pglpm-authoryear,autopunct=false,sorting=ny,sortcites=false,natbib=false,maxcitenames=1,maxbibnames=8,minbibnames=8,giveninits=true,uniquename=false,uniquelist=false,maxalphanames=1,block=space,hyperref=true,defernumbers=false,useprefix=true,sortupper=false,language=british,parentracker=false]{biblatex}
\DeclareSortingScheme{ny}{\sort{\field{sortname}\field{author}\field{editor}}\sort{\field{year}}}
\iffalse\makeatletter%%% replace parenthesis with brackets
\newrobustcmd*{\parentexttrack}[1]{%
  \begingroup
  \blx@blxinit
  \blx@setsfcodes
  \blx@bibopenparen#1\blx@bibcloseparen
  \endgroup}
\AtEveryCite{%
  \let\parentext=\parentexttrack%
  \let\bibopenparen=\bibopenbracket%
  \let\bibcloseparen=\bibclosebracket}
\makeatother\fi
\DefineBibliographyExtras{british}{\def\finalandcomma{\addcomma}}
\renewcommand*{\finalnamedelim}{\addcomma\space}
\setcounter{biburlnumpenalty}{1}
\setcounter{biburlucpenalty}{0}
\setcounter{biburllcpenalty}{1}
\DeclareDelimFormat{multicitedelim}{\addsemicolon\space}
\DeclareDelimFormat{compcitedelim}{\addsemicolon\space}
\DeclareDelimFormat{postnotedelim}{\space}
\ifarxiv\else\addbibresource{../portamanabib.bib}\fi
\renewcommand{\bibfont}{\footnotesize}
%\appto{\citesetup}{\footnotesize}% smaller font for citations
\defbibheading{bibliography}[\bibname]{\section*{#1}\addcontentsline{toc}{section}{#1}%\markboth{#1}{#1}
}
\newcommand*{\citep}{\parencites}
\newcommand*{\citey}{\parencites*}
%\renewcommand*{\cite}{\parencite}
\renewcommand*{\cites}{\parencites}
\providecommand{\href}[2]{#2}
\providecommand{\eprint}[2]{\texttt{\href{#1}{#2}}}
\newcommand*{\amp}{\&}
% \newcommand*{\citein}[2][]{\textnormal{\textcite[#1]{#2}}%\addtocategory{extras}{#2}
% }
\newcommand*{\citein}[2][]{\textnormal{\textcite[#1]{#2}}%\addtocategory{extras}{#2}
}
\newcommand*{\citebi}[2][]{\textcite[#1]{#2}%\addtocategory{extras}{#2}
}
\newcommand*{\subtitleproc}[1]{}
\newcommand*{\chapb}{ch.}

% \def\arxivp{}
% \def\mparcp{}
% \def\philscip{}
% \def\biorxivp{}
% \newcommand*{\arxivsi}{\texttt{arXiv} eprints available at \url{http://arxiv.org/}.\\}
% \newcommand*{\mparcsi}{\texttt{mp\_arc} eprints available at \url{http://www.ma.utexas.edu/mp_arc/}.\\}
% \newcommand*{\philscisi}{\texttt{philsci} eprints available at \url{http://philsci-archive.pitt.edu/}.\\}
% \newcommand*{\biorxivsi}{\texttt{bioRxiv} eprints available at \url{http://biorxiv.org/}.\\}
\newcommand*{\arxiveprint}[1]{%\global\def\arxivp{\arxivsi}%\citeauthor{0arxivcite}\addtocategory{ifarchcit}{0arxivcite}%eprint
\texttt{\urlalt{https://arxiv.org/abs/#1}{arXiv:\hspace{0pt}#1}}%
%\texttt{\href{http://arxiv.org/abs/#1}{\protect\url{arXiv:#1}}}%
%\renewcommand{\arxivnote}{\texttt{arXiv} eprints available at \url{http://arxiv.org/}.}
}
\newcommand*{\mparceprint}[1]{%\global\def\mparcp{\mparcsi}%\citeauthor{0mparccite}\addtocategory{ifarchcit}{0mparccite}%eprint
\texttt{\urlalt{http://www.ma.utexas.edu/mp_arc-bin/mpa?yn=#1}{mp\_arc:\hspace{0pt}#1}}%
%\texttt{\href{http://www.ma.utexas.edu/mp_arc-bin/mpa?yn=#1}{\protect\url{mp_arc:#1}}}%
%\providecommand{\mparcnote}{\texttt{mp_arc} eprints available at \url{http://www.ma.utexas.edu/mp_arc/}.}
}
\newcommand*{\philscieprint}[1]{%\global\def\philscip{\philscisi}%\citeauthor{0philscicite}\addtocategory{ifarchcit}{0philscicite}%eprint
\texttt{\urlalt{http://philsci-archive.pitt.edu/archive/#1}{PhilSci:\hspace{0pt}#1}}%
%\texttt{\href{http://philsci-archive.pitt.edu/archive/#1}{\protect\url{PhilSci:#1}}}%
%\providecommand{\mparcnote}{\texttt{philsci} eprints available at \url{http://philsci-archive.pitt.edu/}.}
}
\newcommand*{\biorxiveprint}[1]{%\global\def\biorxivp{\biorxivsi}%\citeauthor{0arxivcite}\addtocategory{ifarchcit}{0arxivcite}%eprint
\texttt{\urlalt{https://doi.org/10.1101/#1}{bioRxiv doi:\hspace{0pt}10.1101/#1}}%
%\texttt{\href{http://arxiv.org/abs/#1}{\protect\url{arXiv:#1}}}%
%\renewcommand{\arxivnote}{\texttt{arXiv} eprints available at \url{http://arxiv.org/}.}
}
\newcommand*{\osfeprint}[1]{%
\texttt{\urlalt{https://doi.org/10.17605/osf.io/#1}{Open Science Framework doi:10.17605/osf.io/#1}}%
}

\usepackage{graphicx}
%\usepackage{wrapfig}
%\usepackage{tikz-cd}

\PassOptionsToPackage{hyphens}{url}\usepackage[hypertexnames=false]{hyperref}
\usepackage[depth=4]{bookmark}
\hypersetup{colorlinks=true,bookmarksnumbered,pdfborder={0 0 0.25},citebordercolor={0.2667 0.4667 0.6667},citecolor=mypurpleblue,linkbordercolor={0.6667 0.2 0.4667},linkcolor=myredpurple,urlbordercolor={0.1333 0.5333 0.2},urlcolor=mygreen,breaklinks=true,pdftitle={\pdftitle},pdfauthor={\pdfauthor}}
% \usepackage[vertfit=local]{breakurl}% only for arXiv
\providecommand*{\urlalt}{\href}

%%% Layout. I do not know on which kind of paper the reader will print the
%%% paper on (A4? letter? one-sided? double-sided?). So I choose A5, which
%%% provides a good layout for reading on screen and save paper if printed
%%% two pages per sheet. Average length line is 66 characters and page
%%% numbers are centred.
\ifafour\setstocksize{297mm}{210mm}%{*}% A4
\else\setstocksize{210mm}{5.5in}%{*}% 210x139.7
\fi
\settrimmedsize{\stockheight}{\stockwidth}{*}
\setlxvchars[\normalfont] %313.3632pt for a 66-characters line
\setxlvchars[\normalfont]
\setlength{\trimtop}{0pt}
\setlength{\trimedge}{\stockwidth}
\addtolength{\trimedge}{-\paperwidth}
% The length of the normalsize alphabet is 133.05988pt - 10 pt = 26.1408pc
% The length of the normalsize alphabet is 159.6719pt - 12pt = 30.3586pc
% Bringhurst gives 32pc as boundary optimal with 69 ch per line
% The length of the normalsize alphabet is 191.60612pt - 14pt = 35.8634pc
\ifafour\settypeblocksize{*}{32pc}{1.618} % A4
%\setulmargins{*}{*}{1.667}%gives 5/3 margins % 2 or 1.667
\else\settypeblocksize{*}{26pc}{1.618}% nearer to a 66-line newpx and preserves GR
\fi
\setulmargins{*}{*}{1}%gives equal margins
\setlrmargins{*}{*}{*}
\setheadfoot{\onelineskip}{2.5\onelineskip}
\setheaderspaces{*}{2\onelineskip}{*}
\setmarginnotes{2ex}{10mm}{0pt}
\checkandfixthelayout[nearest]
\fixpdflayout
%%% End layout
%% this fixes missing white spaces
\pdfmapline{+dummy-space <dummy-space.pfb}\pdfinterwordspaceon%

%%% Sectioning
\newcommand*{\asudedication}[1]{%
{\par\centering\textit{#1}\par}}
\newenvironment{acknowledgements}{\section*{Thanks}\addcontentsline{toc}{section}{Thanks}}{\par}
\makeatletter\renewcommand{\appendix}{\par
  \bigskip{\centering
   \interlinepenalty \@M
   \normalfont
   \printchaptertitle{\sffamily\appendixpagename}\par}
  \setcounter{section}{0}%
  \gdef\@chapapp{\appendixname}%
  \gdef\thesection{\@Alph\c@section}%
  \anappendixtrue}\makeatother
\counterwithout{section}{chapter}
\setsecnumformat{\upshape\csname the#1\endcsname\quad}
\setsecheadstyle{\large\bfseries\sffamily%
\raggedright}
\setsubsecheadstyle{\bfseries\sffamily%
\raggedright}
%\setbeforesecskip{-1.5ex plus 1ex minus .2ex}% plus 1ex minus .2ex}
%\setaftersecskip{1.3ex plus .2ex }% plus 1ex minus .2ex}
%\setsubsubsecheadstyle{\bfseries\sffamily\slshape\raggedright}
%\setbeforesubsecskip{1.25ex plus 1ex minus .2ex }% plus 1ex minus .2ex}
%\setaftersubsecskip{-1em}%{-0.5ex plus .2ex}% plus 1ex minus .2ex}
\setsubsecindent{0pt}%0ex plus 1ex minus .2ex}
\setparaheadstyle{\bfseries\sffamily%
\raggedright}
\setcounter{secnumdepth}{2}
\setlength{\headwidth}{\textwidth}
\newcommand{\addchap}[1]{\chapter*[#1]{#1}\addcontentsline{toc}{chapter}{#1}}
\newcommand{\addsec}[1]{\section*{#1}\addcontentsline{toc}{section}{#1}}
\newcommand{\addsubsec}[1]{\subsection*{#1}\addcontentsline{toc}{subsection}{#1}}
\newcommand{\addpara}[1]{\paragraph*{#1.}\addcontentsline{toc}{subsubsection}{#1}}
\newcommand{\addparap}[1]{\paragraph*{#1}\addcontentsline{toc}{subsubsection}{#1}}

% Headers and footers
\copypagestyle{manaart}{plain}
\makeheadrule{manaart}{\headwidth}{0.5\normalrulethickness}
\makeoddhead{manaart}{%
{\footnotesize%\sffamily%
\scshape\headauthor}}{}{{\footnotesize\sffamily%
\headtitle}}
\makeoddfoot{manaart}{}{\thepage}{}
\newcommand*\autanet{\includegraphics[height=\heightof{M}]{../autanet.pdf}}
\definecolor{mygray}{gray}{0.333}
\iftypodisclaim%
\ifafour\newcommand\addprintnote{\begin{picture}(0,0)%
\put(245,149){\makebox(0,0){\rotatebox{90}{\tiny\color{mygray}\textsf{This
            document is designed for screen reading and
            two-up printing on A4 or Letter paper}}}}%
\end{picture}}% A4
\else\newcommand\addprintnote{\begin{picture}(0,0)%
\put(176,112){\makebox(0,0){\rotatebox{90}{\tiny\color{mygray}\textsf{This
            document is designed for screen reading and
            two-up printing on A4 or Letter paper}}}}%
\end{picture}}\fi%afourtrue
\makeoddfoot{plain}{}{\makebox[0pt]{\thepage}\addprintnote}{}
\else
\makeoddfoot{plain}{}{\makebox[0pt]{\thepage}}{}
\fi%typodisclaimtrue
\makeoddhead{plain}{}{}{\footnotesize\reporthead}

% \copypagestyle{manainitial}{plain}
% \makeheadrule{manainitial}{\headwidth}{0.5\normalrulethickness}
% \makeoddhead{manainitial}{%
% \footnotesize\sffamily%
% \scshape\headauthor}{}{\footnotesize\sffamily%
% \headtitle}
% \makeoddfoot{manaart}{}{\thepage}{}

\pagestyle{manaart}

\setlength{\droptitle}{-3.9\onelineskip}
\pretitle{\begin{center}\LARGE\sffamily%
\bfseries}
\posttitle{\bigskip\end{center}}

\makeatletter\newcommand*{\atf}{\includegraphics[%trim=1pt 1pt 0pt 0pt,
totalheight=\heightof{@}]{../atblack.png}}\makeatother
\providecommand{\affiliation}[1]{\textsl{\textsf{\footnotesize #1}}}
\providecommand{\epost}[1]{\texttt{\footnotesize\textless#1\textgreater}}
\providecommand{\email}[2]{\href{mailto:#1ZZ@#2 ((remove ZZ))}{#1\protect\atf#2}}

\preauthor{\vspace{-0.5\baselineskip}\begin{center}
\normalsize\sffamily%
\lineskip  0.5em}
\postauthor{\par\end{center}}
\predate{\DTMsetdatestyle{mydate}\begin{center}\footnotesize}
\postdate{\end{center}\vspace{-\medskipamount}}
\usepackage[british]{datetime2}
\DTMnewdatestyle{mydate}%
{% definitions
\renewcommand*{\DTMdisplaydate}[4]{%
\number##3\ \DTMenglishmonthname{##2} ##1}%
\renewcommand*{\DTMDisplaydate}{\DTMdisplaydate}%
}
\DTMsetdatestyle{mydate}


\setfloatadjustment{figure}{\footnotesize}
\captiondelim{\quad}
\captionnamefont{\footnotesize\sffamily%
}
\captiontitlefont{\footnotesize}
\firmlists*
\midsloppy

% handling orphan/widow lines, memman.pdf
% \clubpenalty=10000
% \widowpenalty=10000
% \raggedbottom
% Downes, memman.pdf
\clubpenalty=9996
\widowpenalty=9999
\brokenpenalty=4991
\predisplaypenalty=10000
\postdisplaypenalty=1549
\displaywidowpenalty=1602

\selectlanguage{british}\frenchspacing
%%%%%%%%%%%%%%%%%%%%%%%%%%%%%%%%%%%%%%%%%%%%%%%%%%%%%%%%%%%%%%%%%%%%%%%%%%%%
%%%%%%%%%%%%%%%%%%%%%%%%%%%%%%%%%%%%%%%%%%%%%%%%%%%%%%%%%%%%%%%%%%%%%%%%%%%%
%%%% Paper's details %%%%
\title{\propertitle}
\author{%
\iftrue%
\hspace*{\stretch{1}}%
\parbox{0.5\linewidth}%\makebox[0pt][c]%
{\protect\centering Andr\'es\\%
\footnotesize\epost{\email{ivan.a.davidovich}{ntnu.no}}}%
\hspace*{\stretch{1}}%
\parbox{0.5\linewidth}%\makebox[0pt][c]%
{\protect\centering Luca\\%
\footnotesize\epost{\email{piero.mana}{ntnu.no}}}%
\hspace*{\stretch{1}}%
\else%
Luca\\{\footnotesize\epost{\email{piero.mana}{ntnu.no}}%
%\quad\href{https://orcid.org/0000-0002-6070-0784}{\protect\includegraphics[scale=0.16]{orcid_32x32.png}\textsc{orcid}:0000-0002-6070-0784}%
}\fi%
}

\date{Draft of \today\ (first drafted \firstdraft)}
%\date{\firstpublished; updated \updated}

%@@@@@@@@@@ new macros @@@@@@@@@@
% Common ones - uncomment as needed
%\providecommand{\nequiv}{\not\equiv}
%\providecommand{\coloneqq}{\mathrel{\mathop:}=}
%\providecommand{\eqqcolon}{=\mathrel{\mathop:}}
%\providecommand{\varprod}{\prod}
\newcommand*{\de}{\partialup}%partial diff
\newcommand*{\pu}{\piup}%constant pi
\newcommand*{\delt}{\deltaup}%Kronecker, Dirac
%\newcommand*{\eps}{\varepsilonup}%Levi-Civita, Heaviside
%\newcommand*{\riem}{\zetaup}%Riemann zeta
%\providecommand{\degree}{\textdegree}% degree
%\newcommand*{\celsius}{\textcelsius}% degree Celsius
%\newcommand*{\micro}{\textmu}% degree Celsius
%\newcommand*{\I}{\mathrm{i}}%imaginary unit
%\newcommand*{\e}{\mathrm{e}}%Neper
\newcommand*{\di}{\mathrm{d}}%differential
%\newcommand*{\Di}{\mathrm{D}}%capital differential
%\newcommand*{\planckc}{\hslash}
%\newcommand*{\avogn}{N_{\textrm{A}}}
%\newcommand*{\NN}{\bm{\mathrm{N}}}
%\newcommand*{\ZZ}{\bm{\mathrm{Z}}}
%\newcommand*{\QQ}{\bm{\mathrm{Q}}}
\newcommand*{\RR}{\bm{\mathrm{R}}}
\newcommand*{\CC}{\bm{\mathrm{C}}}
%\newcommand*{\nabl}{\bm{\nabla}}%nabla
%\DeclareMathOperator{\lb}{lb}%base 2 log
%\DeclareMathOperator{\tr}{tr}%trace
%\DeclareMathOperator{\card}{card}%cardinality
%\DeclareMathOperator{\im}{Im}%im part
%\DeclareMathOperator{\re}{Re}%re part
%\DeclareMathOperator{\sgn}{sgn}%signum
%\DeclareMathOperator{\ent}{ent}%integer less or equal to
%\DeclareMathOperator{\Ord}{O}%same order as
%\DeclareMathOperator{\ord}{o}%lower order than
%\newcommand*{\incr}{\triangle}%finite increment
\newcommand*{\defd}{\coloneqq}
\newcommand*{\defs}{\eqqcolon}
%\newcommand*{\Land}{\bigwedge}
%\newcommand*{\Lor}{\bigvee}
%\newcommand*{\lland}{\mathbin{\ \land\ }}
%\newcommand*{\llor}{\mathbin{\ \lor\ }}
%\newcommand*{\lonlyif}{\mathbin{\Rightarrow}}%implies
%\newcommand*{\limplies}{\mathbin{\Rightarrow}}%implies
\newcommand*{\mimplies}{\Rightarrow}%implies
%\newcommand*{\liff}{\mathbin{\Leftrightarrow}}%if and only if
%\newcommand*{\cond}{\mathpunct{|}}%conditional sign (in probabilities)
%\newcommand*{\lcond}{\mathpunct{|\ }}%conditional sign (in probabilities)
%\newcommand*{\bigcond}{\mathpunct{\big|}}%conditional sign (in probabilities)
%\newcommand*{\lbigcond}{\mathpunct{\big|\ }}%conditional sign (in probabilities)
\newcommand*{\suchthat}{\mid}%{\mathpunct{|}}%such that (eg in sets)
%\newcommand*{\bigst}{\mathpunct{\big|}}%such that (eg in sets)
%\newcommand*{\with}{\colon}%with (list of indices)
%\newcommand*{\mul}{\times}%multiplication
%\newcommand*{\inn}{\cdot}%inner product
%\newcommand*{\dotv}{\mathord{\,\cdot\,}}%variable place
%\newcommand*{\comp}{\circ}%composition of functions
%\newcommand*{\con}{\mathbin{:}}%scal prod of tensors
%\newcommand*{\equi}{\sim}%equivalent to 
\renewcommand*{\asymp}{\simeq}%equivalent to 
%\newcommand*{\corr}{\mathrel{\hat{=}}}%corresponds to
%\providecommand{\varparallel}{\ensuremath{\mathbin{/\mkern-7mu/}}}%parallel (tentative symbol)
\renewcommand{\le}{\leqslant}%less or equal
\renewcommand{\ge}{\geqslant}%greater or equal
%\DeclarePairedDelimiter\clcl{[}{]}
\DeclarePairedDelimiter\clop{[}{[}
%\DeclarePairedDelimiter\opcl{]}{]}
%\DeclarePairedDelimiter\opop{]}{[}
%\DeclarePairedDelimiter\abs{\lvert}{\rvert}
%\DeclarePairedDelimiter\norm{\lVert}{\rVert}
\DeclarePairedDelimiter\set{\{}{\}}
%\DeclareMathOperator{\pr}{P}%probability
\newcommand*{\pf}{\mathrm{p}}%probability
\newcommand*{\p}{\mathrm{P}}%probability
%\newcommand*{\tf}{\mathrm{T}}%probability
\renewcommand*{\|}{\mathpunct{|}}
%\newcommand*{\lcond}{\mathpunct{|\ }}%conditional sign (in probabilities)
\newcommand*{\bigcond}{\mathpunct{\big|\ }}%conditional sign (in probabilities)
%\newcommand*{\lbigcond}{\mathpunct{\big|\ }}%conditional sign (in probabilities)
%\newcommand*{\+}{\lor}
%\renewcommand{\*}{\land}
\newcommand*{\sect}{\S}% Sect.~
\newcommand*{\sects}{\S\S}% Sect.~
\newcommand*{\chap}{ch.}%
\newcommand*{\chaps}{chs}%
\newcommand*{\bref}{ref.}%
\newcommand*{\brefs}{refs}%
%\newcommand*{\fn}{fn}%
\newcommand*{\eqn}{eq.}%
\newcommand*{\eqns}{eqs}%
\newcommand*{\fig}{fig.}%
\newcommand*{\figs}{figs}%
\newcommand*{\vs}{{vs}}
%\newcommand*{\etc}{{etc.}}
%\newcommand*{\ie}{{i.e.}}
%\newcommand*{\ca}{{c.}}
%\newcommand*{\eg}{{e.g.}}
\newcommand*{\foll}{{ff.}}
%\newcommand*{\viz}{{viz}}
\newcommand*{\cf}{{cf.}}
%\newcommand*{\Cf}{{Cf.}}
%\newcommand*{\vd}{{v.}}
\newcommand*{\etal}{{et al.}}
%\newcommand*{\etsim}{{et sim.}}
%\newcommand*{\ibid}{{ibid.}}
%\newcommand*{\sic}{{sic}}
%\newcommand*{\id}{\mathte{I}}%id matrix
%\newcommand*{\nbd}{\nobreakdash}%
%\newcommand*{\bd}{\hspace{0pt}}%
%\def\hy{-\penalty0\hskip0pt\relax}
%\newcommand*{\labelbis}[1]{\tag*{(\ref{#1})$_\text{r}$}}
%\newcommand*{\mathbox}[2][.8]{\parbox[t]{#1\columnwidth}{#2}}
%\newcommand*{\zerob}[1]{\makebox[0pt][l]{#1}}
\newcommand*{\tprod}{\mathop{\textstyle\prod}\nolimits}
\newcommand*{\tsum}{\mathop{\textstyle\sum}\nolimits}
%\newcommand*{\tint}{\begingroup\textstyle\int\endgroup\nolimits}
%\newcommand*{\tland}{\mathop{\textstyle\bigwedge}\nolimits}
%\newcommand*{\tlor}{\mathop{\textstyle\bigvee}\nolimits}
%\newcommand*{\sprod}{\mathop{\textstyle\prod}}
%\newcommand*{\ssum}{\mathop{\textstyle\sum}}
%\newcommand*{\sint}{\begingroup\textstyle\int\endgroup}
%\newcommand*{\sland}{\mathop{\textstyle\bigwedge}}
%\newcommand*{\slor}{\mathop{\textstyle\bigvee}}
%\newcommand*{\T}{^\intercal}%transpose
%\newcommand*{\E}{\mathrm{E}}
%\DeclarePairedDelimiter\expp{(}{)}
%\newcommand*{\expe}{\E\expp}%round
%\newcommand*{\expeb}{\E\clcl}%square
%%\newcommand*{\QEM}%{\textnormal{$\Box$}}%{\ding{167}}
%\newcommand*{\qem}{\leavevmode\unskip\penalty9999 \hbox{}\nobreak\hfill
%\quad\hbox{\QEM}}

\definecolor{notecolour}{RGB}{68,170,153}
%\newcommand*{\puzzle}{{\fontencoding{U}\fontfamily{fontawesometwo}\selectfont\symbol{225}}}
\newcommand*{\puzzle}{\maltese}
\newcommand{\mynote}[1]{ {\color{notecolour}\puzzle\ #1}}
%\newcommand*{\widebar}[1]{{\mkern1.5mu\skew{2}\overline{\mkern-1.5mu#1\mkern-1.5mu}\mkern 1.5mu}}

%\DeclareMathOperator*{\argsup}{arg\,sup}

\newcommand{\explanation}[4][t]{%\setlength{\tabcolsep}{-1ex}
%\smash{
\begin{tabular}[#1]{c}#2\\[0.5\jot]\rule{1pt}{#3}\\#4\end{tabular}}%}
\newcommand*{\ptext}[1]{\text{\small #1}}
%@@@@ Custom macros for this file @@@@
\newcommand*{\dob}{degree of belief}
\newcommand*{\dobs}{degrees of belief}
\newcommand*{\yI}{\varIota}
\newcommand*{\yD}{D}
\newcommand*{\yt}{\bm{t}}
\newcommand*{\yf}{\bm{f}}
\newcommand*{\xo}[1]{x^{(#1)}}
\newcommand*{\yo}[1]{y^{(#1)}}
\newcommand*{\yT}{T'}
\newcommand*{\yq}{\bm{q}}
\newcommand*{\yH}{D}
%@@@@@@@@@@ new macros end @@@@@@@@@@

\firmlists
\begin{document}
\captiondelim{\quad}\captionnamefont{\footnotesize}\captiontitlefont{\footnotesize}
\selectlanguage{british}\frenchspacing

%%% Title and abstract %%%
\maketitle
\ifpublic
\abstractrunin
\abslabeldelim{}
\renewcommand*{\abstractname}{}
\setlength{\absleftindent}{0pt}
\setlength{\absrightindent}{0pt}
\setlength{\abstitleskip}{-\absparindent}
\begin{abstract}\labelsep 0pt%
  \noindent The use of generalization error for testing and the phenomenon
  of overtraining are explained geometrically and given a Bayesian
  interpretation.
% \par%\\[\jot]
% \noindent
% {\footnotesize PACS: ***}\qquad%
% {\footnotesize MSC: ***}%
%\qquad{\footnotesize Keywords: ***}
\end{abstract}\fi

\selectlanguage{british}\frenchspacing
% \asudedication{\small ***}
% \vspace{\bigskipamount}

% \setlength{\epigraphwidth}{.7\columnwidth}
% %\epigraphposition{flushright}
% \epigraphtextposition{flushright}
% %\epigraphsourceposition{flushright}
% \epigraphfontsize{\footnotesize}
% \setlength{\epigraphrule}{0pt}
% %\setlength{\beforeepigraphskip}{0pt}
% %\setlength{\afterepigraphskip}{0pt}
% \epigraph{\emph{text}}{source}

\iftrue
\noindent\emph{\footnotesize Note: Dear Reader \amp\ Peer, this manuscript is
  being peer-reviewed by you. Thank you.}
\fi

\section{Generalization tests, overtraining, probability calculus}
\label{sec:gentest_overtr}

In neural-network learning, and more generally in machine learning, it's
customary to choose among candidate neural nets by testing them against a
\emph{generalization} data set separate from the \emph{training} data set
used to find the network parameters. Let's call this a \emph{generalization
  test}; it may assume different forms, including cross-validation. The
candidate neural nets differ in their architecture -- number of nodes,
weights, non-linear functions -- or in the error function and weight
regularizer used to train them.

The generalization test can also be used during the learning phase: it is
sometimes observed that the generalization error reaches a minimum before
the learning phase -- that is, the minimization of the training error -- is
complete. This phenomenon is called \emph{overtraining}. In this case the
learning phase is stopped when the generalization error reaches a minimum.

MacKay \parentext{\cite*{mackay1992,mackay1992b}; see also
  \cite[]{bishop2006}}, combining the analyses of Tishby, Levin, Solla
\citep{tishbyetal1989,levinetal1990} and Gull \citey{gull1989}, gives an
insightful interpretation of machine and neural-nets learning from the
point of view of the probability calculus. For a more thorough discussion
of the statements in the next two paragraphs, see MacKay
\citey{mackay1992,mackay1992b} and Bishop \citey[esp. \sects~1.2--3,
3.3--4, 5.7]{bishop2006}.

From the point of view of the probability calculus it makes sense to choose
among candidates by comparing their plausibilities conditional on the
training data set:
$\pf(\ptext{candidate} \| \ptext{data}, \ptext{initial info})$. If we
initially judge the candidates to be equally plausible, their
plausibilities conditional on the data are proportional to the plausibility
of the data conditional on the candidates, that is, to the likelihoods of
the candidates:
$\pf(\ptext{data} \| \ptext{candidate}, \ptext{initial info})$, also called
the \emph{evidence} for the candidate.

The evidence and the generalization test often lead to the same choice of
candidate, and the evidence has several advantages compared to the
generalization test. Yet, they aren't equivalent \citep[for a thorough
discussion of their differences see][esp. \sect~6.6]{mackay1992}[esp.
\sect~4.1]{mackay1992b}. In particular it may happen that a candidate with
poor evidence yields a good generalization test. This often indicates that
the whole set of candidates represents our belief about the particular data
poorly. So, even from the point of view of the probability calculus, the
generalization test can be very useful.

\bigskip

The generalization test is not derived by an application of the probability
calculus, however, unlike the evidence. Its idea is nevertheless intuitive and
sound.

In this note we show that the generalization test can in fact be derived
from the rules of probability. The derivation also shows why this test can
indicate that the whole sets of candidates is a poor representation of our
beliefs about the data.

\section{Geometric preliminaries}
\label{sec:geom_prelim}

\section{Learning models and extreme models}
\label{sec:learning_extreme}

Central to our Bayesian understanding of the generalization test is the
difference between a \emph{learning} plausibility model and a non-learning
plausibility model, also called \emph{extremal} or independent, and the
relationship between the two. This relationship, extensively studied by
Lauritzen
\citey{lauritzen1974,lauritzen1974_r1974,lauritzen1982_r1988,lauritzen1984},
is closely connected with the notion of sufficient statistics and
exponential families of distributions [refs***].



\mynote{to be modified}
Consider sets of input quantities $\yt\defd(t_1, t_2, \dotsc)$ and
corresponding output quantities $(x_1, x_2, \dotsc)$.
% In the following the input data are often omitted from the conditional of
% the plausibilities for brevity.
A partially exchangeable plausibility model is a \dob\ about the output
data, conditional on the input data, that is invariant under every relabelling of
those output data that have the same input values. For example, 
\begin{multline}
  \label{eq:example_partial_exch}
  \pf(x_1=1,x_2=2,x_3=3,x_4=4, \dotsc \|
  t_1=5,t_2=6,t_3=5,t_4=6, \dotsc \, \yI)
  ={}\\
  \pf(x_1=3,x_2=2,x_3=1,x_4=4, \dotsc \|
  t_1=5,t_2=6,t_3=5,t_4=6,\dotsc \, \yI)
  ={}\\
  \pf(x_1=3,x_2=4,x_3=1,x_4=2, \dotsc \|
  t_1=5,t_2=6,t_3=5,t_4=6,\dotsc \, \yI),
\end{multline}
that is, we have the same \dob\ about $x_1=1,x_2=2,x_3=3,x_4=4$ and about
$x_1=3,x_2=2,x_3=1,x_4=4$, where the values of $x_1$ and $x_3$ have been
exchanged, because these correspond to input data having the same value,
$t_1=t_3=5$. This \dob\ also equals the one about
$x_1=3,x_2=4,x_3=1,x_4=2$, where the values of $x_2$ and $x_4$ have also
been exchanged, because these correspond to input data having the same
value, $t_2=t_4=6$. And so on. However, this \dob\ can be different from
the one about $x_1=2,x_2=1,x_3=3,x_4=4$, where the values of $x_1$ and
$x_2$ have been exchanged, because these correspond to different input
values: $t_1 \ne t_2$.


A partially exchangeable model is completely determined if we specify our
\dob\ about the relative frequencies of $x$ values for each distinct input
value $t$, in an indefinitely long sequence. [continue explanation***] This
\dob\ is therefore represented by a density $\pf(\set{\yf_t} \| \yI)$ on a
space equivalent to a product of simplices, one for each different value
that $t$ can assume. When this density is positive over the whole space the
model is called \emph{non-parametric}. This space acquires a huge
dimensionality as the numbers of possible $t$ and $x$ values increase. For
this reason we often specify singular densities, which are positive only
over a lower-dimensional subset of this space. The corresponding model is
called \emph{parametric}, because we introduce parameters as coordinates on
the lower-dimensional subset.



When the model is updated or trained with data $\yD$, the corresponding
density $\pf(\set{\yf_t} \| \yI)$ is updated to
$\pf(\set{\yf_t} \| \yD, \yI)$, which yields different predictions. This is
an example of a \emph{learning} model. As the number of training data
increases the density concentrates on a point of the lower-dimensional subset.


Generally, we specify this density 







\section{Bayesian interpretation of the generalization test}
\label{sec:bayes_gentest}

Consider sets of input data $\yt\defd(t_1, t_2, \dotsc)$ and corresponding output
data $(x_1, x_2, \dotsc)$. % In the following the input data are often
% omitted from the conditional of the plausibilities for brevity.

The evidence is our \dob\ about the output data, given the input data and
our initial information and assumptions $\yI$:
\begin{equation}
  \label{eq:evidence_def}
  \pf(x_1, x_2, \dotsc \| t_1, t_2, \dotsc,\, \yI).
\end{equation}
It's important to note that our joint \dob\ about the outputs is not the
product of our marginal \dobs:
\begin{multline}
  \label{eq:evidence_not_product}
  \pf(x_1, x_2, x_3 \dotsc \| \yt,\, \yI)\\
  \!\begin{aligned}[b]
      &=\pf(x_1 \| x_2, x_3, \dotsc,\, \yt,\, \yI) \times
        \pf(x_2 \| x_3, \dotsc,\, \yt,\, \yI) \times 
        \pf(x_3 \| \dotsc,\, \yt,\, \yI) \times \dotsb
  \\
  &\ne
    \pf(x_1 \| \yt,\, \yI) \times
    \pf(x_2 \| \yt,\, \yI) \times
    \pf(x_3\| \yt,\, \yI) \times \dotsb
\end{aligned}
\end{multline}






\clearpage
\hrule

*** Old draft ***

\section{Synopsis}
\label{sec:synopsis}

The literature usually associates overtraining with overfitting and
too large number of parameters.

The relation between these aspects is more subtle than that. In fact,
overtraining can be a consequence of \emph{underfitting}. We have a clearer
view approaching this question from the point of view of probability
theory.

Consider a classification problem: given input
$x \in \set{1,\dotsc, M}$, we want to predict output
$y \in \set{1,\dotsc,N}$. Let's consider both discrete.

Two cases must be distinguished: whether every input has several possible
outputs, or just one. The second case is a special instance of the first,
and could be simply treated as the prediction of a \enquote{function}
$x\mapsto y$. But both cases are instances of inference with a
\emph{partially exchangeable} model.


Our assumptions are summarized in the proposition $\yI$. We assume that the
probability of $y$, for each kind of input $x$, is exchangeable. Thus the
probability for several pairs
\begin{equation}
  \label{eq:prob_several_pairs}
  \pf(\yo{T},\dotsc, \yo{1} \|\xo{T}, \dotsc, \xo{1}, \yI)
\end{equation}
must have a partially exchangeable form [refs]: for every kind of input
$x$ we consider the $N$ relative frequencies
$\yf_x \defd (f_{1\|x},\dotsc, f_{N\|x})$ of the $N$ possible outputs. The
probability above is given by this integral:
\begin{multline}
  \label{eq:prob_several_pairs_exch}
  \pf(\yo{T},\dotsc, \yo{1} \|\xo{T}, \dotsc, \xo{1}, \yI)
  ={}\\
  \idotsint
  \biggl[ \prod_{x=1}^M \prod_t^{\xo{t}=x} f_{\yo{t}\|x} \biggr]
  \;
  \pf(\yf_1,\dotsc,\yf_M \| \yI)
  \;\di\yf_1 \dotsm \di\yf_M.
\end{multline}

The complicated look of this formula doesn't do justice to its simple and
intuitive interpretation:

Choose a kind of input $x$. We're judging that the probability for a very
large sequence of outputs generated with this input doesn't depend on their
order. Suppose we knew the relative frequencies of each kind of output in
this sequence, $(f_{1\|x},\dotsc,f_{N\|x}) \defs \yf_x$, and knew nothing
else. Then out of symmetry reasons we should give a probability $f_{y\|x}$
of observing outcome $y$ at the next observation. This is just a
\enquote{drawing without replacement} problem. At the subsequent
observations with the same input we give again the same probabilities to
each $y$, because the sequence is so large that we approximate
\enquote{drawing without replacement} with \enquote{drawing with
  replacement}. The probability for a sequence of outputs
$\yo{1},\dotsc,\yo{T}$ from the same input $x$ is then
\begin{equation}
  \pf(\yo{T}, \dotsc,\yo{1}\|\text{\footnotesize same input $x$},\yf_x, \yI) =
  f_{\yo{T}\|x}\times\dotsm\times f_{\yo{1}\|x}.
  \label{eq:input_m_indep_output}
\end{equation}
Within a larger sequence of outputs from all possible inputs, the outputs
$\yo(t)$ coming from input $x$ are those for which $\xo(t)=x$. Thus we can
write the formula above more generally as
\begin{equation}
 \prod_t^{\xo{t}=x} f_{\yo{t}\|x}
  \label{eq:input_m_indep_output_label_x}
\end{equation}
The above reasoning applies for each kind of input $x$. Thus the
probability for a sequence of outputs coming from different inputs is the
product of the probabilities above for all different $x$:
\begin{equation}
  \pf(\text{\footnotesize all outputs} \|\text{\footnotesize inputs}, \yf_1,\dotsc,\yf_M, \yI) =
\prod_{x=1}^M \prod_t^{\xo{t}=x} f_{\yo{t}\|x}.
  \label{eq:input_m_indep_output_all_x}
\end{equation}


Now what if we don't know the relative frequencies $\yf_x$, for any input?
Then we assign a probability distribution over their possible values% ,
% $\pf(\yf_1,\dotsc,\yf_M \| \yI)\,\di\yf_1\dotsm\di\yf_M$,
and use the law of total probability:
\begin{multline}
  \label{eq:freqs_total_prob}
  \pf(\text{\footnotesize all outputs} \|\text{\footnotesize inputs}, \yI) ={}\\
  \idotsint \pf(\text{\footnotesize all outputs} \|\text{\footnotesize inputs},
  \yf_1,\dotsc,\yf_M, \yI)
  \; \pf(\yf_1,\dotsc,\yf_M \| \yI)\,\di\yf_1,\dotsc,\di\yf_M.
\end{multline}
Substituting the explicit expression~\eqref{eq:input_m_indep_output_all_x}
in this formula we obtain formula~\eqref{eq:freqs_total_prob}. In summary,
\begin{equation}
  \label{eq:prob_several_pairs_exch_explanation}
  \textcolor{mygreen}{\idotsint}
  \biggl[ \textcolor{mypurpleblue}{\explanation{$\displaystyle\prod_{x=1}^M$}{3em}{\clap{\parbox{10em}{\centering\footnotesize product over all inputs}}}} \textcolor{myred}{\explanation{$\smash{\displaystyle\prod_t^{\xo{t}=x} f_{\yo{t}\|x}}$}{0.67em}{\clap{\quad\parbox[t]{8em}{\centering\footnotesize probability for outputs\\from same input}}}} \biggr]
  \;
  \textcolor{mygreen}{\explanation{$\pf(\yf_1,\dotsc,\yf_M \| \yI)
    \,\di\yf_1 \dotsm \di\yf_M$}{3em}{\clap{\parbox{12em}{\centering\footnotesize uncertainty over all frequencies}}}}.
\end{equation}


The possible frequencies give one input,
$\yf_x \equiv (f_{1\|x}, \dotsc, f_{N\|x})$, belong to the
$(N-1)$-dimensional simplex
\begin{equation}
  \label{eq:simplex_N}
  \varDelta \defd \set[\big]{(f_1, \dotsc,f_N) \suchthat f_i \ge 0, \tsum_i f_i =1},
\end{equation}
and the collection of possible frequencies $(\yf_1,\dotsc, \yf_M)$ belongs
to the $M$-fold Cartesian product $\varDelta^M$. From now on we denote
$\yf \defd (\yf_1,\dotsc,\yf_M)$.

\bigskip

The probability for a new sequence of $\yT$ outputs given their inputs and
given that we've learned a previous sequence of $T$ input-output pairs is
determined by Bayes's theorem:
\begin{equation}
  \label{eq:prob_new_several_pairs}
  \begin{gathered}
  \pf(\yo{\yT},\dotsc, \yo{T+1} \|\xo{\yT}, \dotsc, \xo{T+1},\;
  \yo{T},\xo{T},\dotsc, \yo{1},\xo{1},\; \yI) ={}\\
  \frac{\displaystyle
    \int
  \biggl[ \prod_{x=1}^M\, \prod_{t=T+1,\dotsc,\yT}^{\xo{t}=x} f_{\yo{t}\|x} \biggr]
  \;  \pf(\yf \| \yI)  \;\di\yf
}{\displaystyle
\int
  \biggl[ \prod_{x=1}^M\, \prod_{t=1,\dotsc,T}^{\xo{t}=x} f_{\yo{t}\|x} \biggr]
  \;  \pf(\yf \| \yI)  \;\di\yf
}.
\end{gathered}
\end{equation}
This formula is equivalent to~\eqref{eq:prob_several_pairs_exch} with and
updated distribution for the frequencies:
\begin{equation}
  \label{eq:update_pr_freqs}
  \begin{gathered}
  \pf(\yf \| \yo{T},\xo{T},\dotsc, \yo{1},\xo{1},\; \yI) \;\di\yf={}\\
  \frac{
  \bigl[ \prod_{x=1}^M\, \prod_{t=1,\dotsc,T}^{\xo{t}=x} f_{\yo{t}\|x} \bigr]
  \;  \pf(\yf \| \yI) 
}{
  \int  \bigl[ \prod_{x=1}^M\, \prod_{t=1,\dotsc,T}^{\xo{t}=x} f_{\yo{t}\|x} \bigr]
  \;  \pf(\yf \| \yI) \;\di\yf
  } \;\di\yf.
\end{gathered}
\end{equation}

The form of this updated distribution has important consequences for our
learning process.

If the number of learned data is enough large compared with the numbers
$M$, $N$ of possible inputs and outputs and with the magnitude of the
initial distribution for the frequencies, and if the latter is strictly
positive, then the updated distribution becomes very peaked on the
collection of relative frequencies $(\yq_1,\dotsc,\yq_M)$ of the learned
outputs for all input values. This can be seen from the asymptotic
expression in terms of the relative entropy (Kullback-Leibler divergence)
$\yH$,
  \begin{equation}
    \label{eq:asympt_pf_KL}
    \pf(\yf \| \yo{T},\xo{T},\dotsc, \yo{1},\xo{1},\; \yI)
    \propto
  \exp\bigl[ -\tsum_x T_x\yH(\yq_x \| \yf_x)  \bigr]\;
    \pf(\yf \| \yI),
  \end{equation}
where $T_x$ is the number of observations with input $x$, with $\sum_x T_x=T$.

If the number of learned data is small compared with the dimensions of the
input and output spaces, then the initial distribution for the frequencies
$\pf(\yf \| \yI)\,\di\yf$ greatly influence our
inference~\eqref{eq:prob_new_several_pairs}. This distribution determines
two important ***



\section{Utility functions and probabilities}
\label{sec:utilities_probs}

Utility of behaving as if proposition $B \in X$ is true given that proposition
$A \in X$ is true: $c(B \|A)$. Probability for $A$ given $D,\yI$: $\p(A \|
D, \yI)$. Optimal decision is $B$ that maximizes
\begin{equation}
  \label{eq:optimal_decision}
  \sum_A c(B \| A)\, \p(A \|D,\yI).
\end{equation}

Now consider different probabilities for all $A$ given the same data $D$:
$\p(A\|D, \yI')$. The decision $B$ will still be the same if we use a new
utility
\begin{equation}
  \label{eq:new_utility}
  c'(B \| A) \defd c(B \|A)\,\frac{\p(A \|D,\yI)}{\p(A \|D,\yI')}.
\end{equation}
So the same choice can be made with a different probability, if the utility
is appropriately changed, provided $p(A \|D,\yI')>0$ for all $A$.

This leads to a slightly more general view than Tishby \etal's
\citey{tishbyetal1989,levinetal1990} and Mackay's \citey{mackay1992,mackay1992b}


\newpage

\textcolor{white}{If you find this you can claim a postcard from me.}
%

\setlength{\intextsep}{0.5ex}% with wrapfigure
%\begin{figure}[p!]%{r}{0.4\linewidth} % with wrapfigure
%  \centering\includegraphics[trim={12ex 0 18ex 0},clip,width=\linewidth]{maxent_saddle.png}\\
%\caption{***}\label{fig:comparison_a5}
%\end{figure}% exp_family_maxent.nb


\ifpublic
\begin{acknowledgements}
  \ldots to Mari \amp\ Miri for continuous encouragement and affection, and
  to Buster Keaton and Saitama for filling life with awe and inspiration.
  To the developers and maintainers of \LaTeX, Emacs, AUC\TeX, Open Science
  Framework, Python, Inkscape, Sci-Hub for making a free and unfiltered
  scientific exchange possible.
%\rotatebox{15}{P}\rotatebox{5}{I}\rotatebox{-10}{P}\rotatebox{10}{\reflectbox{P}}\rotatebox{-5}{O}.
\sourceatright{\autanet}
\end{acknowledgements}
\fi

\renewcommand*{\appendixpagename}{Appendix}
\renewcommand*{\appendixname}{Appendix}
%\appendixpage
%\appendix

%%%%%%%%%%%%%%% BIB %%%%%%%%%%%%%%%

\defbibnote{prenote}{{\footnotesize (\enquote{de $X$} is listed under D,
    \enquote{van $X$} under V, and so on, regardless of national
    conventions.)\par}}
% \defbibnote{postnote}{\par\medskip\noindent{\footnotesize% Note:
%     \arxivp \mparcp \philscip \biorxivp}}

\printbibliography[prenote=prenote%,postnote=postnote
]


\end{document}
---------- cut text ----------------


%%% Local Variables: 
%%% mode: LaTeX
%%% TeX-PDF-mode: t
%%% TeX-master: t
%%% End: 
